\setchapterpreamble[u]{\margintoc}
\chapter{Introduzione}
\labch{Introduzione}

In questa tesi esploreremo gli Euler Tour Trees 
(ETT), una struttura dati dinamica introdotta da 
Henzinger e King \cite{HenzingerKing1999}, 
progettata per mantenere informazioni su foreste 
soggette a modifiche strutturali.

L'obiettivo è supportare operazioni di: 
\begin{itemize}
    \item \textbf{link(u,v)}: aggiungere un arco $(u,v)$ unendo due alberi.
    \item \textbf{cut(u,v)}: rimuovere l'arco $(u,v)$ dividendo l'albero in due.
    \item \textbf{isConnected(u,v)}: verificare se $u$ e $v$ appartengono allo stesso albero.
\end{itemize}

Il presente lavoro di tesi si propone di sistematizzare ed estendere il framework teorico relativo 
alla struttura dati introdotta da Henzinger e King \cite{HenzingerKing1999}. L’obiettivo principale è 
stato quello di colmare le lacune informative presenti nel contributo originale, descrivendole in maniera 
esaustiva e rigorosa di tutti i dettagli implementativi necessari alla piena comprensione del suo funzionamento.\\

Un elemento di forte originalità della ricerca risiede nella scelta dell'architettura sottostante: 
a differenza delle implementazioni convenzionali basate sugli Splay Tree, si è scelto di adottare i Red-Black Tree. 
Tale approccio garantisce il mantenimento di un costo computazionale logaritmico nel caso peggiore, 
offrendo una maggiore stabilità prestazionale rispetto alle strutture auto-bilancianti basate sull'accesso.

